\subsection{The unordered fixed-point equivalence}
\label{sec:p03-abiteboul-vianu-completion}

The construction orders tuple types, not the original universe.
Its ordered structure retains the information needed by a fixed
formula without choosing an order on the original elements.

The type-invariant method of Dawar, Lindell and Weinstein is presented
in \cite[Chapter~11]{Libkin2004}. The evaluation below also retains
nested fixed-point binders and their parameters.

\begin{hbconvention}{def:p03-av-setting}{Finite structures and tuple types}
Use the finite vocabularies of
\cref{def:p03-w03-p03-056-vocabulary}: relation symbols of positive
arity and possibly constants, with finite nonempty universes.
For \(k\geq1\), two tuples \(\bar a,\bar b\in A^k\)
have the same \(k\)-variable type when they satisfy the same
first-order formulas using the variables \(x_1,\ldots,x_k\).
This is \(\equiv^k_{\omega\omega}\) of
\cref{def:p03-w03-p03-056-equivalence}, with the variables renamed.
Write \(\bar a[i\leftarrow u]\) for replacement of coordinate
\(i\). This is notation for a tuple, not a function symbol added
to the object language.
\end{hbconvention}

\begin{hblemma}{thm:p03-av-type-order}{proved=annex:p03.av.1}{Ordering the tuple types}
For each vocabulary and \(k\), there is an LFP formula
\(D_k(\bar x,\bar y)\) such that, in every finite structure,
\[
 E_k(\bar x,\bar y)
 \quad\defeq\quad
 \neg D_k(\bar x,\bar y)\land\neg D_k(\bar y,\bar x)
\]
defines equality of \(k\)-variable types, and \(D_k\) induces a
strict linear order on the \(E_k\)-classes. The formula is uniform
in the input structure. It need not define a linear order on \(A\).
\end{hblemma}

\begin{hbdefinition}{def:p03-av-quotient}{The ordered type structure}
Fix the formula \(D_k\) just constructed. The universe of
\(\operatorname{Typ}_k(\mathbf A)\) is \(A^k/E_k\). Its vocabulary
is finite and depends only on \(k\) and the original vocabulary.
It contains the induced strict order and a constant \(c_\ast\)
naming its least element, together with the following relations.
\begin{enumerate}
\item For each atomic formula \(\alpha\) in the fixed variables
      and constants, a unary predicate \(U_\alpha\), true of
      \([\bar a]\) exactly when \(\mathbf A\models\alpha[\bar a]\).
\item For each coordinate \(i\), a binary relation \(S_i\), where
      \(S_i([\bar a],[\bar b])\) means that some representatives
      agree outside coordinate \(i\). Equivalently,
      \(\exists u\,E_k(\bar a[i\leftarrow u],\bar b)\).
\item Each coordinate map \(\theta\) selects, at every position,
      one of the \(k\) fixed variables or an original constant.
      Include a binary relation \(P_\theta\) with graph
      \[
       P_\theta([\bar a],[\bar b])
             \quad\Longleftrightarrow\quad E_k(\bar a_\theta,\bar b).
      \]
      The \(i\)-th entry of \(\bar a_\theta\) is the value of
      \(\theta(i)\) under \(\bar a\).
\end{enumerate}
There are finitely many atomic formulas and maps in these lists.
The next lemma verifies that the definitions do not depend on
representatives.
\end{hbdefinition}

\begin{hblemma}{thm:p03-av-pullback}{proved=annex:p03.av.2}{Interpreting the ordered quotient}
\(\operatorname{Typ}_k(\mathbf A)\) is well defined and has between
one and \(|A|^k\) elements. Its relations and its quotient equality
have uniform LFP definitions on representative tuples. Every LFP
sentence \(\eta\) in its vocabulary has an LFP sentence
\(\eta^{\leftarrow}\) in the original vocabulary satisfying
\[
 \operatorname{Typ}_k(\mathbf A)\models\eta
 \quad\Longleftrightarrow\quad
 \mathbf A\models\eta^{\leftarrow}.
\]
The construction handles \(c_\ast\) by its definable equivalence
class; it does not choose a definable representative from that class.
\end{hblemma}

\begin{hblemma}{thm:p03-av-quotient-space}{proved=annex:p03.av.3}{Evaluation on types}
For every PFP sentence \(\varphi\) there are a fixed \(k\)
and an isomorphism-invariant polynomial-space decision problem
\(Q_\varphi\) on finite ordered structures of the vocabulary
of \(\operatorname{Typ}_k\), with its distinguished constant, such that
\[
 \mathbf A\models\varphi
 \quad\Longleftrightarrow\quad
 \operatorname{Typ}_k(\mathbf A)\in Q_\varphi.
\]
Space is polynomial in the size of the ordered type structure,
not in the size of \(A\). The bound \(k\) may increase under
renaming and nesting; no sharp bound by the original number of
printed variable names is asserted. Parameters and free relation
variables of inner fixed-point matrices are retained during evaluation.
\end{hblemma}

Under \(\mathsf P=\mathsf{PSPACE}\), ordered capture supplies an
LFP sentence for \(Q_\varphi\), and the pullback lemma removes
the auxiliary ordered structure. This establishes the implication from complexity-class equality in
\cref{thm:p03-w03-p03-059-abiteboul-vianu}; its annex proof also
passes from sentences to queries with their original free variables.
It does not assert that LFP captures all polynomial-time queries on
unordered structures: only the specified PFP query is recovered from
its tuple-type information.
